

\documentclass[11pt,a4paper]{article}
\usepackage[utf8]{inputenc}
\usepackage[T1]{fontenc}
\usepackage[french]{babel}
\usepackage{lmodern}
\usepackage{microtype}
\usepackage[a4paper,margin=1.65cm,headheight=15pt]{geometry}
\usepackage{amsmath,amssymb,mathtools}
\usepackage{array,tabularx,multicol,booktabs}
\usepackage{enumitem}
\usepackage[most]{tcolorbox}
\usepackage{xcolor}
\usepackage{tikz}
\usetikzlibrary{arrows.meta,calc,positioning}
\usepackage{pdflscape}
\usepackage{fancyhdr}
\usepackage{hyperref}
\usepackage{pagecolor}
\usepackage{needspace}

\definecolor{fondpage}{RGB}{247,248,250}
\definecolor{encre}{RGB}{31,35,40}
\definecolor{bleu}{RGB}{40,91,135}
\definecolor{vert}{RGB}{55,125,92}
\definecolor{orange}{RGB}{178,105,42}
\definecolor{violet}{RGB}{101,77,145}
\definecolor{rouge}{RGB}{170,72,72}
\definecolor{gris}{RGB}{86,91,99}
\definecolor{bleufonce}{RGB}{28,55,120}
\definecolor{bleuclair}{RGB}{238,243,255}
\definecolor{grisclair}{RGB}{245,245,245}

\hypersetup{colorlinks=true,linkcolor=bleufonce,urlcolor=bleufonce,
pdftitle={Parcours de revision -- Fonction exponentielle -- Premiere specialite},pdfauthor={}}

\setlength{\parindent}{0pt}
\setlength{\parskip}{0.25em}
\renewcommand{\arraystretch}{1.13}
\setlength{\tabcolsep}{4pt}
\setlist[itemize]{leftmargin=1.25em,itemsep=0.15em,topsep=0.2em}
\setlist[enumerate,1]{label=\textbf{\arabic*.}, leftmargin=1.05cm, itemsep=0.35em, topsep=0.2em}
\setlist[enumerate,2]{label=\textbf{\alph*.}, leftmargin=0.85cm, itemsep=0.25em, topsep=0.15em}

\pagestyle{fancy}
\fancyhf{}
\cfoot{\thepage}

\newcommand{\R}{\mathbb{R}}
\newcommand{\N}{\mathbb{N}}
\newcommand{\sourceq}[1]{ {\scriptsize\textcolor{bleufonce}{(site Premiere q#1)}}}
\newcounter{question}
\newcounter{reponse}

\newenvironment{blocquestions}[1]{%
  \begin{tcolorbox}[enhanced,breakable,colback=bleuclair,colframe=bleufonce,boxrule=0.6pt,arc=2mm,left=2mm,right=2mm,top=2mm,bottom=1.5mm,title={\bfseries #1},coltitle=white,colbacktitle=bleufonce,before skip=3mm,after skip=3mm, fontupper=\large]
  \large
  \begin{enumerate}[leftmargin=*,label=\textbf{\arabic*.},itemsep=0.82em,topsep=0.5em]
  \setcounter{enumi}{\value{question}}
}{%
  \setcounter{question}{\value{enumi}}
  \end{enumerate}
  \end{tcolorbox}
}

\newenvironment{blocreponses}[1]{%
  \begin{tcolorbox}[enhanced,breakable,colback=bleuclair,colframe=bleufonce,boxrule=0.6pt,arc=2mm,left=2mm,right=2mm,top=2mm,bottom=1.5mm,title={\bfseries #1},coltitle=white,colbacktitle=bleufonce,before skip=3mm,after skip=3mm, fontupper=\large]
  \large
  \begin{enumerate}[leftmargin=*,label=\textbf{\arabic*.},itemsep=0.42em,topsep=0.5em]
  \setcounter{enumi}{\value{reponse}}
}{%
  \setcounter{reponse}{\value{enumi}}
  \end{enumerate}
  \end{tcolorbox}
}

\newtcolorbox{correctionbox}[1]{enhanced,breakable,colback=grisclair,colframe=black!55,boxrule=0.55pt,arc=2mm,left=2mm,right=2mm,top=2mm,bottom=1.5mm,title=\textbf{#1},coltitle=black,colbacktitle=black!12,before skip=2mm,after skip=2mm}
\newtcolorbox{methodebox}[1]{enhanced,breakable,colback=white,colframe=bleu!70!black,boxrule=0.55pt,arc=2mm,left=2mm,right=2mm,top=2mm,bottom=1.5mm,title=\textbf{#1},coltitle=white,colbacktitle=bleu!70!black,before skip=2mm,after skip=2mm}
\newtcolorbox{exoBox}[2]{enhanced,breakable,colback=white,colframe=#1!70!black,boxrule=0.55pt,arc=2mm,left=2mm,right=2mm,top=2mm,bottom=1.5mm,title=\textbf{#2},coltitle=white,colbacktitle=#1!70!black,before skip=3mm,after skip=3mm}

\newcommand{\titreSujet}[2]{%
  \subsection*{#1}
  \addcontentsline{toc}{subsection}{#1}
  \textit{Référence / inspiration : #2}\par\medskip\hrule\medskip
}
\newcommand{\qcm}[4]{%
\begin{center}\begin{tabularx}{0.96\linewidth}{@{}XX@{}}
\textbf{a.} #1 & \textbf{b.} #2\\[1mm]
\textbf{c.} #3 & \textbf{d.} #4
\end{tabularx}\end{center}}

\tcbset{enhanced,arc=2mm,outer arc=2mm,boxrule=0.42pt,left=1.15mm,right=1.15mm,top=0.75mm,bottom=0.75mm,boxsep=0.35mm,before skip=0.8mm,after skip=0.8mm,fonttitle=\bfseries\footnotesize,coltitle=encre,before upper=\raggedright\fontsize{7.15}{7.65}\selectfont}
\newtcolorbox{fichebox}[3][]{colback=white,colframe=#2!72!black,colbacktitle=#2!18,coltitle=#2!50!black,borderline west={0.9mm}{0pt}{#2!78!black},title={#3},drop fuzzy shadow=black!5,#1}
\newcommand{\tagcol}[2]{\begin{tcolorbox}[enhanced,colback=#1!11,colframe=#1!45!black,boxrule=0.42pt,arc=2mm,left=1mm,right=1mm,top=0.45mm,bottom=0.45mm,before skip=0mm,after skip=0.85mm,fontupper=\bfseries\scriptsize,colupper=#1!55!black]\centering #2\end{tcolorbox}}
\newtcolorbox{panel}[1]{colback=fondpage,colframe=fondpage,boxrule=0pt,arc=0mm,left=0mm,right=0mm,top=0mm,bottom=0mm,height=168mm,valign=top,before skip=0mm,after skip=0mm}

\newcommand{\titrepage}{\begin{tcolorbox}[colback=white,colframe=bleu!70!black,boxrule=0.65pt,arc=3mm,left=2.5mm,right=2.5mm,top=1.15mm,bottom=1.15mm,before skip=0mm,after skip=1.4mm,drop fuzzy shadow=black!10]
{\Large\bfseries Parcours de révision -- Fonction exponentielle}\hfill {\normalsize Première spécialité -- sans calculatrice}
\end{tcolorbox}}
\newcommand{\courbeExp}{%
\begin{center}
\begin{tikzpicture}[scale=0.52]
\draw[->,gray] (-2.1,0)--(2.8,0) node[right]{$x$};
\draw[->,gray] (0,-0.2)--(0,4.4) node[above]{$y$};
\foreach \x in {-2,-1,0,1,2}{\draw[gray] (\x,0.06)--(\x,-0.06) node[below] {\scriptsize $\x$};}
\foreach \y/\lab in {1/1,2/2,3/3}{\draw[gray] (0.06,\y)--(-0.06,\y) node[left] {\scriptsize $\lab$};}
\draw[very thick,bleu,domain=-2:1.45,samples=90] plot (\x,{exp(\x)});
\node[bleu] at (1.65,3.5) {$y=e^x$};
\draw[dashed,gray] (1,0)--(1,2.718);
\fill (1,2.718) circle (1.5pt) node[right] {\scriptsize $e$};
\end{tikzpicture}
\end{center}}

\newcommand{\tableVarSimple}[5]{%
\begin{center}
\begin{tikzpicture}[x=1cm,y=0.7cm]
\draw[fill=grisclair] (0,2) rectangle (8,3);
\draw (0,0) rectangle (8,3);
\draw (1.45,0)--(1.45,3);
\draw (0,1)--(8,1); \draw (0,2)--(8,2);
\node at (0.72,2.5) {$x$};
\node at (0.72,1.5) {$f'(x)$};
\node at (0.72,0.5) {$f$};
\node at (1.85,2.5) {$#1$};
\node at (4.65,2.5) {$#2$};
\node at (7.45,2.5) {$#3$};
\node at (2.9,1.5) {$#4$};
\node at (4.65,1.5) {$0$};
\node at (6.4,1.5) {$#5$};
\end{tikzpicture}
\end{center}}

\newcommand{\ficheRecap}{%
\begin{landscape}
\pagecolor{fondpage}
\thispagestyle{empty}
\titrepage
\begin{tabularx}{\linewidth}{@{}X@{\hspace{1.8mm}}X@{\hspace{1.8mm}}X@{\hspace{1.8mm}}X@{}}
\begin{panel}{}
\tagcol{bleu}{Définition et courbe}
\begin{fichebox}{bleu}{Définition}
La fonction exponentielle est l'unique fonction $\exp$ dérivable sur $\mathbb R$ telle que :
\[
\exp(0)=1 \quad\text{et}\quad \exp'=\exp.
\]
On note :
\[
\exp(x)=e^x.
\]
\end{fichebox}
\begin{fichebox}{bleu}{Valeurs de référence}
\[
e^0=1,\qquad e^1=e.
\]
Le nombre $e$ est environ égal à $2{,}718$.
\end{fichebox}
\begin{fichebox}{bleu}{Courbe et signe}
\courbeExp
Pour tout réel $x$ :
\[
e^x>0.
\]
\end{fichebox}
\end{panel}
&
\begin{panel}{}
\tagcol{vert}{Propriétés algébriques}
\begin{fichebox}{vert}{Formules}
Pour tous réels $a$ et $b$ :
\[
e^{a+b}=e^a e^b,
\qquad
 e^{a-b}=\frac{e^a}{e^b}.
\]
\[
(e^a)^n=e^{na},
\qquad
 e^{-a}=\frac1{e^a}.
\]
\end{fichebox}
\begin{fichebox}{vert}{Attention}
On ne développe pas $e^{a+b}$ :
\[
e^{a+b}\neq e^a+e^b.
\]
\end{fichebox}
\end{panel}
&
\begin{panel}{}
\tagcol{orange}{Dériver et varier}
\begin{fichebox}{orange}{Dérivées de référence}
\[
(e^x)'=e^x.
\]
Pour tous réels $a$ et $b$ :
\[
\bigl(e^{ax+b}\bigr)'=a e^{ax+b}.
\]
\end{fichebox}
\begin{fichebox}{orange}{Produit avec exponentielle}
Pour dériver un produit, on utilise :
\[
(uv)'=u'v+uv'.
\]
Exemple :
\[
\bigl(xe^x\bigr)'=e^x+xe^x=(x+1)e^x.
\]
\end{fichebox}
\end{panel}
&
\begin{panel}{}
\tagcol{violet}{Équations et inéquations}
\begin{fichebox}{violet}{Croissance}
La fonction exponentielle est strictement croissante sur $\mathbb R$.\\[0.3em]
Donc :
\[
e^a=e^b \Longleftrightarrow a=b.
\]
\[
e^a<e^b \Longleftrightarrow a<b.
\]
\end{fichebox}
\begin{fichebox}{violet}{Méthode}
Pour résoudre une équation du type :
\[
e^{2x-1}=e^{x+3},
\]
on compare les exposants :
\[
2x-1=x+3.
\]
\end{fichebox}
\end{panel}
\end{tabularx}
\clearpage
\nopagecolor
\end{landscape}}

\begin{document}
\begin{center}
{\LARGE\bfseries Parcours de révision -- Fonction exponentielle}\\[1mm]
{\large Première générale -- spécialité mathématiques}\\[1mm]
{\normalsize Sans calculatrice}
\end{center}
\tableofcontents\newpage
\phantomsection\addcontentsline{toc}{section}{Fiche récapitulative}
\ficheRecap

\section{Questions flash}
\setcounter{question}{0}
\begin{blocquestions}{Questions flash 1 à 12}
\item Calculer $e^0$.
\item Calculer $e^1$.
\item Simplifier $e^{x+2}e^{3}$.
\item Simplifier $\dfrac{e^{x+4}}{e^x}$.
\item Simplifier $e^{2x}e^{1-x}$.
\item Simplifier $\dfrac{e^{3x}}{e^x}$.
\item Simplifier $(e^x)^4$.
\item Simplifier $e^{-x}e^{2x}$.
\item Dériver $f(x)=e^x$.
\item Dériver $f(x)=3e^x$.
\item Dériver $f(x)=e^{2x}$.
\item Dériver $f(x)=e^{-x}$.
\end{blocquestions}

\begin{blocquestions}{Questions flash 13 à 24}
\item Dériver $f(x)=e^{3x-1}$.
\item Dériver $f(x)=2e^{x}-x$.
\item Dériver $f(x)=xe^x$.
\item Dériver $f(x)=(x-2)e^x$.
\item Résoudre $e^x=1$.
\item Résoudre $e^{x+2}=e^5$.
\item Résoudre $e^{2x}=e^{x+4}$.
\item Résoudre $e^{3x-1}=e^{x+5}$.
\item Résoudre $e^{x}<e^4$.
\item Résoudre $e^{2x-1}\ge e^3$.
\item Donner le signe de $e^x$.
\item Donner le sens de variation de $e^x$.
\end{blocquestions}

\begin{blocquestions}{Questions flash 25 à 36}
\item Donner le signe de $(x-3)e^x$.
\item Donner le signe de $-(x+1)e^x$.
\item Factoriser $e^{2x}-e^x$.
\item Résoudre $e^x(e^x-1)=0$.
\item Dériver $f(x)=x+e^{1-x}$.
\item Dériver $f(x)=(2x+1)e^x$.
\item Si $f'(x)=(x-1)e^x$, donner le signe de $f'(x)$.
\item Si $f'(x)=-2e^{x}$, donner le sens de variation de $f$.
\item Calculer $e^2/e^{-1}$.
\item Simplifier $e^{x-2}\times e^{2-x}$.
\item Résoudre $e^{4-x}=e^{1+x}$.
\item Résoudre $e^{x+1}\le e^{2x-3}$.
\end{blocquestions}


\section{QCM d'automatismes}
Pour chaque question, une seule réponse est correcte.
\begin{enumerate}
\item $e^{a+b}=$\par\vspace{0.8mm}\qcm{$e^a+e^b$}{$e^ae^b$}{$e^{ab}$}{$a e^b$}
\item $e^{-a}=$\par\vspace{0.8mm}\qcm{$-e^a$}{$\frac1{e^a}$}{$e^a$}{$0$}
\item La dérivée de $e^{3x}$ est\par\vspace{0.8mm}\qcm{$e^{3x}$}{$3e^{3x}$}{$3e^x$}{$e^{x+3}$}
\item L'équation $e^{x-1}=e^4$ a pour solution\par\vspace{0.8mm}\qcm{$3$}{$4$}{$5$}{$e^5$}
\item Pour tout réel $x$, $e^x$ est\par\vspace{0.8mm}\qcm{négatif}{nul}{positif}{parfois nul}
\item Le signe de $(x+2)e^x$ est le signe de\par\vspace{0.8mm}\qcm{$x+2$}{$e^x$}{$x$}{$x-2$}
\item $e^{2x}/e^{x-1}=$\par\vspace{0.8mm}\qcm{$e^{x+1}$}{$e^{x-1}$}{$e^{3x-1}$}{$e^{x}$}
\item La fonction exponentielle est\par\vspace{0.8mm}\qcm{décroissante}{constante}{croissante}{paire}
\end{enumerate}

\section{Sujets type bac}


\titreSujet{Sujet 1 -- Comparaison de deux fonctions exponentielles}{inspiré d’exercices E3C sur équations, inéquations et position relative de courbes}
On considère les fonctions $f$ et $g$ définies sur $\R$ par :
\[
f(x)=e^{2x}+e^2
\qquad\text{et}\qquad
g(x)=(1+e^2)e^x.
\]
On note $\mathcal C_f$ et $\mathcal C_g$ leurs courbes représentatives.
\begin{enumerate}
\item Montrer que, pour tout réel $x$ :
\[
f(x)-g(x)=e^{2x}-(1+e^2)e^x+e^2.
\]
\item Factoriser le trinôme :
\[
X^2-(1+e^2)X+e^2.
\]
\item En posant $X=e^x$, en déduire que, pour tout réel $x$ :
\[
f(x)-g(x)=(e^x-1)(e^x-e^2).
\]
\item Déterminer les abscisses des points d'intersection de $\mathcal C_f$ et $\mathcal C_g$.
\item Donner les coordonnées exactes de ces points d'intersection.
\item Étudier le signe de $f(x)-g(x)$ selon les valeurs de $x$.
\item En déduire la position relative de $\mathcal C_f$ et $\mathcal C_g$.
\item Résoudre l'inéquation $f(x)\le g(x)$.
\item Sans calculer de valeurs décimales, comparer $2e^2$ et $e+e^3$.
\end{enumerate}

\titreSujet{Sujet 2 -- Étude d'une fonction avec exponentielle}{inspiré d’exercices E3C sur dérivation, variations et tangente avec exponentielle}
On considère la fonction $f$ définie sur $[-2;4]$ par :
\[
f(x)=(x-2)e^x.
\]
\begin{enumerate}
\item Calculer $f(0)$.
\item Montrer que, pour tout réel $x$ :
\[
f'(x)=(x-1)e^x.
\]
\item Étudier le signe de $f'(x)$ sur $[-2;4]$.
\item Dresser le tableau de variations de $f$ sur $[-2;4]$.
\item Déterminer la valeur exacte du minimum de $f$ sur $[-2;4]$.
\item Étudier le signe de $f(x)$ sur $[-2;4]$.
\item Déterminer l'équation de la tangente à la courbe de $f$ au point d'abscisse $0$.
\end{enumerate}


\titreSujet{Sujet 3 -- Modèle en contexte}{inspiré d’exercices E3C de modélisation par une fonction exponentielle}
La concentration d'un produit dans le sang, en unité arbitraire, est modélisée pendant les cinq premières heures par la fonction $C$ définie sur $[0;5]$ par :
\[
C(t)=t e^{1-t}.
\]
\begin{enumerate}
\item Calculer $C(0)$ et $C(1)$.
\item Montrer que, pour tout $t\in[0;5]$ :
\[
C'(t)=(1-t)e^{1-t}.
\]
\item Étudier le signe de $C'(t)$ sur $[0;5]$.
\item Dresser le tableau de variations de $C$ sur $[0;5]$.
\item À quel instant la concentration est-elle maximale ? Quelle est cette concentration maximale ?
\item Calculer $C(2)$.
\item On admet que $e<3$. Montrer que $C(2)>\dfrac23$.
\item Interpréter ce résultat dans le contexte.
\end{enumerate}


\end{document}
